\chaptermark{Clausura cardinal bidireccional del continuo analítico HMT}
\label{chap:83-19}

Este capítulo no propone un generador alternativo del continuo. Su premisa
causal es el bloque único ya construido
\[
 \begin{aligned}
 \mathrm{APP}&\longrightarrow\mathrm{TRIT}\longrightarrow\mathrm{TPK}
 \longrightarrow\text{estado enriquecido único}\\
 &\longrightarrow\text{estado común, árbol y medida genealógica}\\
 &\longrightarrow\text{operaciones intrínsecas correlativas}\\
 &\longrightarrow\text{ensamblaje natural, terminalidad y límite}
 \longrightarrow\mathcal C_{\rm cont}^{\rm disc}.
 \end{aligned}
\]
Las cinco construcciones consustanciales desarrolladas en el capítulo
anterior han emergido correlativamente sobre ese mismo estado antes de la
publicación real, sin pretender agotar todas sus propiedades. Lo que sigue es una presentación
posterior del cociente, de su memoria temporal y de su potencia; no selecciona
retrospectivamente ninguna de esas construcciones.

\section{La célula TRIT completa y la extensión temporal bidireccional}

El signo temporal no se separa del estado que orienta. La célula local de
partida es
\begin{equation}
 c=(i,j;R^+,Q^+,R^\times,Q^\times;\tau),
 \qquad \tau\in\{-1,0,+1\},
 \label{p12-83-c19-s1:hmtch:eq-celula-trit-completa}
\end{equation}
donde
\[
 R^+=\operatorname{dr}_9(i+j),\qquad
 Q^+=\frac{i+j-R^+}{9},
\]
\[
 R^\times=\operatorname{dr}_9(ij),\qquad
 Q^\times=\frac{ij-R^\times}{9}.
\]
Los residuos constituyen la lectura visible y los cocientes conservan la
memoria necesaria para invertir el transporte. El trit temporal actúa sobre
esta célula entera:
\begin{equation}
\begin{aligned}
 +1&\longmapsto \text{retorno, memoria retenida y pasado},\\
 0&\longmapsto \text{umbral de evaluación y presente},\\
 -1&\longmapsto \text{apertura bidireccional y futuro}.
\end{aligned}
\label{p12-83-c19-s1:hmtch:eq-lector-temporal-trit}
\end{equation}
Así, dos apariciones del mismo signo no son el mismo estado si difieren sus
cocientes, su hoja, su acarreo o su registro genealógico.

Sea \(X\) una carta compacta del límite superviviente APP--TRIT--TPK y sea
\(F:X\to X\) su actualización completa, continua y sobreyectiva. La
sobreyectividad de los
truncamientos y la prolongación superviviente permiten formar la extensión
natural
\begin{equation}
 X_F^{\leftrightarrow}
 :=\left\{(x_k)_{k\in\mathbb Z}\in X^{\mathbb Z}:
 F(x_k)=x_{k+1}\ \text{para todo }k\in\mathbb Z\right\}.
 \label{p12-83-c19-s1:hmtch:eq-extension-natural-two-way}
\end{equation}
El desplazamiento
\[
 \sigma((x_k)_{k\in\mathbb Z})=(x_{k+1})_{k\in\mathbb Z}
\]
ya no pierde la dirección inversa.

\begin{theorem}[Extensión natural two-way del estado enriquecido]
\label{p12-83-c19-s1:hmtch:thm-extension-natural-two-way}
El espacio \(X_F^{\leftrightarrow}\) es compacto y

\[
 \sigma:X_F^{\leftrightarrow}\longrightarrow X_F^{\leftrightarrow}
\]
es un homeomorfismo, con
\[
 \sigma^{-1}((x_k))=(x_{k-1}).
\]
Por tanto, ida y retorno son morfismos continuos inversos sobre el objeto
enriquecido. La aparente discontinuidad sólo aparece si primero se olvidan
las coordenadas que hacen invertible la historia.
\end{theorem}

\begin{proof}
El producto \(X^{\mathbb Z}\) es compacto. La condición
\(F(x_k)=x_{k+1}\) es cerrada porque \(F\) es continua; por tanto
\(X_F^{\leftrightarrow}\) es compacto. El desplazamiento y el
desplazamiento inverso son continuos en la topología producto y preservan
las ecuaciones de compatibilidad. Sus composiciones son la identidad.
\end{proof}

La inversión material de la ruta vive en la completación grupoidal del
transporte. Sea \(C_M=-\operatorname{rev}\) la conjugación de memoria. Su
acción contravariante se expresa por el cuadrado de reversión
\[
 F(x_k)=x_{k+1}
 \quad\Longrightarrow\quad
 F(C_Mx_{k+1})=C_Mx_k.
\]
Sobre la extensión natural se define
\begin{equation}
 (\Theta u)_k:=C_M(x_{-k}),
 \qquad u=(x_k)_{k\in\mathbb Z}.
 \label{p12-83-c19-s1:hmtch:eq-definicion-espejo-two-way}
\end{equation}
La identidad de reversión garantiza que \(\Theta u\) vuelve a satisfacer
las ecuaciones de compatibilidad. Por tanto,
\begin{equation}
 \Theta^2=I,
 \qquad
 \Theta\sigma\Theta=\sigma^{-1}.
 \label{p12-83-c19-s1:hmtch:eq-espejo-inversion-two-way}
\end{equation}
Esta identidad no identifica el grupo ternario con el signo binario: el
TRIT orienta retorno, umbral y apertura, mientras la involución \(C_M\)
implementa la conjugación de las dos orientaciones temporales.

Para \(n\geq0\), denotemos por \(S_n^{\leftrightarrow}\) el conjunto finito
de ventanas temporales admisibles
\((x_{-n},\ldots,x_0,\ldots,x_n)\), con cada estado truncado además a
profundidad genealógica \(n\). Las restricciones simultáneas de tiempo y
profundidad forman un
sistema inverso y
\begin{equation}
\mathscr O_{\leftrightarrow}^{\rm HMT}
 :=\varprojlim_n S_n^{\leftrightarrow}
 \cong X_F^{\leftrightarrow}.
 \label{p12-83-c19-s1:hmtch:eq-biordinal-genealogico}
\end{equation}
Cada historia posee soporte temporal de tipo
\(\omega^*+1+\omega\), pero lleva además la célula
\eqref{p12-83-c19-s1:hmtch:eq-celula-trit-completa}, la ruta y el registro genealógico. La aplicación
que envía una órbita a todas sus ventanas centradas es un homeomorfismo y
entrelaza \(\sigma\) con la sucesión genealógica. Éste es el biordinal
genealógico: no el ordinal de von Neumann obtenido tras olvidar las
etiquetas.

El objeto anterior es una carta compacta de la parte fraccionaria. El objeto
global que publica toda la recta es el atlas numerable
\begin{equation}
 \mathscr O_{\mathbb R,\leftrightarrow}^{\rm HMT}
 :=\mathbb Z\times\mathscr O_{\leftrightarrow}^{\rm HMT}.
 \label{p12-83-c19-s1:hmtch:eq-biordinal-global}
\end{equation}
La coordenada \(\mathbb Z\) sólo traslada la carta; no borra el registro genealógico
bidireccional.

\section{Calendario unitario, memoria modular y modo estable}

La estructura temporal no se apoya únicamente en una metáfora del retorno.
Sobre
\(\mathcal H_{\rm clk}=\mathbb C^{108}\), sea
\(\zeta=\exp(2\pi i/108)\),
\[
 U_{\rm clk}|j\rangle=|j+1\rangle,
 \qquad
 |\widetilde k\rangle
 =\frac1{\sqrt{108}}\sum_{j=0}^{107}\zeta^{jk}|j\rangle,
\]
y definamos
\begin{equation}
 H_{\rm clk}
 =\hbar_{\rm int}\omega_0
 \sum_{k=0}^{107}k\,|\widetilde k\rangle\langle\widetilde k|,
 \qquad
 \omega_0=\frac{2\pi}{108t_0}.
 \label{p12-83-c19-s1:hmtch:eq-hamiltoniano-reloj}
\end{equation}
Entonces
\begin{equation}
 U_{\rm step}
 =\exp\!\left(-\frac{i}{\hbar_{\rm int}}H_{\rm clk}t_0\right)
 =U_{\rm clk},
 \qquad U_{\rm step}^{108}=I,
 \label{p12-83-c19-s1:hmtch:eq-propagador-reloj}
\end{equation}
y, para
\(Z_{\rm clk}|j\rangle=\zeta^j|j\rangle\),
\[
 |\psi_n\rangle:=U_{\rm step}^n|j_0\rangle,
\]
\begin{equation}
 \langle Z_{\rm clk}\rangle_n
 :=\langle\psi_n|Z_{\rm clk}|\psi_n\rangle
 =\zeta^{j_0+n}
 \label{p12-83-c19-s1:hmtch:eq-respuesta-primitiva-108}
\end{equation}
tiene periodo primitivo \(108\).

En el sector modular de borde
\(\mathcal H_A=(\mathbb C^3)^{\otimes36}\), sea \(\rho_A>0\) precisamente
el estado fiel caracterizado por la descomposición que sigue y sea
\(\Delta_4^{\rm mod}\) el salto modular declarado. El generador que conserva
la procedencia de los dos canales del borde es
\begin{equation}
 K_{\rm HHM}^{(A)}=-\log\rho_A
 =\kappa_A I+\Delta_4^{\rm mod}(90N_{90}+120N_{120}).
 \label{p12-83-c19-s1:hmtch:eq-hamiltoniano-modular-memoria}
\end{equation}
Si
\[
 N=N_{90}+N_{120},
 \qquad
 D=90N_{90}+120N_{120},
\]
entonces
\begin{equation}
 N_{90}=4N-\frac{D}{30},
 \qquad
 N_{120}=\frac{D}{30}-3N.
 \label{p12-83-c19-s1:hmtch:eq-reconstruccion-canales-modulares}
\end{equation}
Así, el observable total y el Hamiltoniano modular reconstruyen
conjuntamente las dos memorias de procedencia. Por otra parte, la dinámica
cilíndrica de Perron conserva el modo
\begin{equation}
 V_n=\varphi^n,
 \qquad
 M_n=\varphi^{-2n}=V_n^{-2},
 \qquad
 M_n=\left(\frac{a_n}{a_0}\right)^{-6}\qquad(D=3).
 \label{p12-83-c19-s1:hmtch:eq-modo-estable-memoria}
\end{equation}

\begin{theorem}[Composición finita exacta del calendario CT108 y la memoria
modular \(90/120\)]
\label{p12-83-c19-s1:hmtch:thm-composicion-calendario-memoria}
Los operadores
\eqref{p12-83-c19-s1:hmtch:eq-hamiltoniano-reloj} y
\eqref{p12-83-c19-s1:hmtch:eq-hamiltoniano-modular-memoria} son autoadjuntos en sus
respectivos espacios. Tras fijar una escala energética \(E_\partial>0\), el
operador del producto
\begin{equation}
 H_{\rm tw}
 :=H_{\rm clk}\otimes I
 +I\otimes E_\partial(K_{\rm HHM}^{(A)}-\kappa_A I)
 \label{p12-83-c19-s1:hmtch:eq-hamiltoniano-two-way-total}
\end{equation}
es autoadjunto y sus dos sumandos conmutan fuertemente. Para estados cuya
marginal de reloj es \(|j_0\rangle\langle j_0|\), el observable
\(Z_{\rm clk}\otimes I\) conserva su respuesta primitiva de periodo \(108\).
Simultáneamente, \(K_{\rm HHM}^{(A)}\), \(N_{90}\) y \(N_{120}\) son
observables conservados del factor modular. El par \((N,D)\), con
\(N=N_{90}+N_{120}\), reconstruye exactamente \((N_{90},N_{120})\) mediante
\eqref{p12-83-c19-s1:hmtch:eq-reconstruccion-canales-modulares}.
\end{theorem}

\begin{proof}
El primer operador posee una descomposición espectral real finita, y el
segundo es el logaritmo de un estado fiel; ambos son autoadjuntos. Actúan en
factores tensoriales distintos y, por tanto, sus resoluciones espectrales
conmutan. La exponencial de su suma factoriza. La primitividad de la respuesta
procede de la primitividad de \(\zeta\); la reconstrucción de canales es la
inversión de la matriz
\(\left(\begin{smallmatrix}1&1\\90&120\end{smallmatrix}\right)\), de
determinante \(30\).
\end{proof}

\begin{corollary}[Modo estable de la memoria cilíndrica]
\label{p12-83-c19-s1:hmtch:cor-modo-estable-memoria}
En el sector de Perron de rutas cerradas,
\(M_n=\varphi^{-2n}=V_n^{-2}\) es el modo espectral exacto y estable dado
por \eqref{p12-83-c19-s1:hmtch:eq-modo-estable-memoria}.
\end{corollary}

El calendario CT108, la memoria modular y el modo de Perron son tres lectores
tipados del mismo antecedente genealógico; no son el mismo operador. Junto con
la extensión natural y el espejo, forman el registro temporal interno
\[
 \mathcal T_{\rm HMT}
 :=\bigl(X_F^{\leftrightarrow},\sigma,\Theta;
 H_{\rm clk},K_{\rm HHM}^{(A)};M_\bullet\bigr).
\]
Éste es el contenido exacto del cierre temporal HMT: dos direcciones
continuas inversas, reloj primitivo, memoria recuperable y modo estable. No
se sustituye el estado TPK completo por el cursor \(108\). Una identificación
adicional de los tres lectores con una única dinámica física perturbativamente
estable requeriría un entrelazador común; esa identificación no se utiliza en
la clausura genealógica ni en la reconstrucción del continuo.

\section{Publicaciones analíticas de historias HMT}

Sea \(H\) un espacio polaco de historias: un límite inverso cerrado de
niveles finitos o un atlas numerable de tales límites. La topología está
generada por cilindros de prefijo. Esta elección conserva la genealogía y
permite tipar con exactitud qué operaciones descriptivas se admiten.

\begin{definition}[Publicación HMT analítica]
\label{p12-83-c19-s1:hmtch:def-publicacion-analitica}
Una publicación analítica es un conjunto \(A\subseteq\mathbb R\) para el
que existen un boreliano
\[
 C\subseteq H\times\mathbb N^{\mathbb N}
\]
y un lector boreliano \(F:H\to\mathbb R\) tales que
\[
 A=\{F(x):\exists z\in\mathbb N^{\mathbb N},\ (x,z)\in C\}.
\]
El testigo \(z\) puede codificar una ruta, una prolongación, una memoria o
una certificación numerable. Los lectores continuos son casos particulares.
\end{definition}

Esta clase contiene los cilindros, los predicados borelianos sobre
historias, las operaciones numerables, las imágenes por lectores continuos
y las publicaciones que afirman la existencia de un testigo real. No
contiene por definición toda la potencia \(\mathcal P(\mathbb R)\): las
uniones no numerables, los selectores arbitrarios, los cocientes no suaves y
la recursión de longitud \(\omega_1\) exigen datos adicionales.

\begin{theorem}[Dicotomía perfecta de las publicaciones HMT]
\label{p12-83-c19-s1:hmtch:thm-dicotomia-perfecta}
Toda publicación HMT analítica \(A\subseteq\mathbb R\) satisface exactamente
una de las dos alternativas:
\[
 |A|\leq\aleph_0
 \qquad\text{o}\qquad
 |A|=2^{\aleph_0}.
\]
En la segunda alternativa existe una inyección continua
\[
 j:2^{\mathbb N}\hookrightarrow A.
\]
Por tanto, toda publicación analítica no numerable contiene una copia del
espacio de Cantor y, en particular, un subconjunto perfecto.
\end{theorem}

\begin{proof}
La proyección existencial de un boreliano de un producto de espacios polacos
es analítica, y una imagen boreliana de un espacio boreliano estándar es de
nuevo analítica. Por tanto, \(A\) es un subconjunto analítico de
\(\mathbb R\).

Para hacer explícita la alternativa no numerable, represéntese \(A\) como
la proyección de un cerrado
\[
 D\subseteq\mathbb R\times\mathbb N^{\mathbb N}.
\]
Para cada palabra finita \(s\in\mathbb N^{<\mathbb N}\), sea
\[
 A_s:=\{x:\exists z\supseteq s,\ (x,z)\in D\}.
\]
Si \(A_s\cap K\) es no numerable para cierto intervalo cerrado \(K\), el
conjunto de sus puntos que no son de condensación es numerable, porque
\(\mathbb R\) posee una base numerable. Se eligen dos puntos de condensación
distintos, dos intervalos cerrados disjuntos dentro de \(K\), y para cada
intervalo una prolongación estricta de \(s\) cuya proyección sigue siendo no
numerable. Las dos prolongaciones no necesitan ser incompatibles: la
separación de las ramas queda garantizada por los intervalos. La identidad
\[
 A_s=\bigcup_{n\in\mathbb N}A_{s^\frown n}
\]
garantiza que, en cada entorno elegido, alguna extensión conserva la no
numerabilidad.

Repitiendo el procedimiento se obtienen, para cada
\(u\in2^{<\mathbb N}\), un intervalo \(K_u\) y un prefijo \(s_u\) tales que
\[
 K_{u0}\cap K_{u1}=\varnothing,
 \qquad
 K_{ui}\subseteq K_u,
 \qquad
 \operatorname{diam}K_u\leq2^{-|u|},
\]
y \(s_{ui}\) prolonga estrictamente a \(s_u\). Para
\(\beta\in2^{\mathbb N}\), los intervalos
\(K_{\beta|n}\) determinan un único punto \(j(\beta)\), y los prefijos
\(s_{\beta|n}\) determinan un testigo \(z_\beta\). La clausura de \(D\)
implica
\[
 (j(\beta),z_\beta)\in D,
\]
de modo que \(j(\beta)\in A\). Los intervalos hermanos disjuntos prueban la
inyectividad, y el control de diámetros prueba la continuidad de \(j\).

Así,
\[
 2^{\aleph_0}\leq |A|\leq|\mathbb R|=2^{\aleph_0},
\]
lo que cierra la segunda alternativa. Si el proceso no parte de un conjunto
no numerable, se obtiene la primera.
\end{proof}

\begin{corollary}[Hipótesis del continuo analítica HMT]
\label{p12-83-c19-s1:hmtch:cor-ch-analitica-hmt}
Sea
\begin{equation}
 \aleph_{1,\mathrm{HMT}}^{\rm an}
 :=\min\left\{|A|:
 A\text{ es una publicación HMT analítica no numerable}\right\}.
 \label{p12-83-c19-s1:hmtch:eq-aleph1-analitico-hmt}
\end{equation}
Este símbolo denota el cardinal inicial relativo a la clase de publicaciones
analíticas HMT; no denota el ordinal extensional \(\omega_1\).
Entonces
\begin{equation}
 \boxed{
 \aleph_{1,\mathrm{HMT}}^{\rm an}
 =2^{\aleph_0}
 =|\mathbb R|
 =\left|\mathscr O_{\mathbb R,\leftrightarrow}^{\rm HMT}/
 \!\sim_{\mathfrak R}\right|.}
 \label{p12-83-c19-s1:hmtch:eq-ch-analitica-hmt}
\end{equation}
Es decir, no existe un cardinal intermedio entre lo numerable y el continuo
dentro de las publicaciones analíticas producidas por lectores HMT.
\end{corollary}

\begin{proof}
El conjunto de ventanas de profundidad finita
\[
 \mathscr O_{<\omega}^{\rm HMT}
 :=\bigcup_{n<\omega}S_n^{\leftrightarrow}
\]
es numerable, pues es unión numerable de conjuntos finitos. La capa completa
sobreyecta sobre \(\mathbb R\) por
\eqref{p12-83-c18-s2:hmtch:eq-homeomorfismo-two-way-real}, y está contenida en un producto
numerable de alfabetos finitos; por tanto tiene cardinal
\(2^{\aleph_0}\). El
\cref{p12-83-c19-s1:hmtch:thm-dicotomia-perfecta} excluye cualquier tamaño analítico
estrictamente intermedio.
\end{proof}

\begin{corollary}[Dicotomía de las fibras]
\label{p12-83-c19-s1:hmtch:cor-dicotomia-fibras}
Sea \(F:H\to Y\) un lector boreliano entre espacios polacos y sea
\(y\in Y\). Si el dominio de historias admitidas es analítico, la fibra
\[
 \{x:F(x)=y\}
\]
es numerable o posee cardinal \(2^{\aleph_0}\).
\end{corollary}

\begin{proof}
La fibra es la intersección del dominio analítico con la preimagen
boreliana del singleton \(\{y\}\); por tanto es analítica. Se aplica el
\cref{p12-83-c19-s1:hmtch:thm-dicotomia-perfecta}.
\end{proof}

\section{Comparación entre el cierre cardinal genealógico y su reducción ordinal extensional}

Sea \(\operatorname{WO}\subseteq2^{\mathbb N}\) el conjunto de códigos
reales de buenos órdenes numerables, permitiendo que el campo sea un
subconjunto de \(\mathbb N\). La aplicación
\[
 \operatorname{ot}:\operatorname{WO}\longrightarrow\omega_1
\]
asigna a cada código su tipo de orden y es sobreyectiva.

\begin{proposition}[Inyección ordinal en la recta]
\label{p12-83-c19-s1:hmtch:prop-omega1-en-reales}
En ZFC existe una inyección
\[
 \iota:\omega_1\hookrightarrow\mathbb R.
\]
Puede levantarse a un subsistema binario del espacio de historias HMT.
\end{proposition}

\begin{proof}
Por elección, la sobreyección \(\operatorname{ot}\) admite una sección
\(s:\omega_1\to\operatorname{WO}\). Equivalentemente, se fija un buen
orden de \(\operatorname{WO}\) y se toma en cada fibra el primer código.
La aplicación de Cantor
\[
 c(a):=\sum_{n\geq0}\frac{2a(n)}{3^{n+1}}
\]
es inyectiva de \(2^{\mathbb N}\) en \([0,1]\). Por tanto
\[
 \iota:=c\circ s
\]
es la inyección buscada. El embebimiento de la cinta binaria en la cinta
ternaria y la inversa Hensel \(\Psi^{-1}\) proporcionan el levantamiento al
espacio genealógico.
\end{proof}

La elección de \(s\) no es natural respecto de los refinamientos ni es
producida por un prefijo finito. Esta proposición muestra que la flecha
\(\omega_1\hookrightarrow\mathbb R\) no es el contenido de la hipótesis
clásica del continuo. La dirección inversa pertenece al ordinal extensional
desnudo. En HMT, en cambio, la flecha natural ya construida es el
homeomorfismo \eqref{p12-83-c18-s2:hmtch:eq-homeomorfismo-two-way-real}, cuyo dominio
conserva las dos direcciones temporales.

\section{Identidad cardinal del continuo analítico HMT: \(\aleph_{1,\mathrm{HMT}}^{\mathrm{an}}=2^{\aleph_0}\) y comparación con la hipótesis clásica del continuo}

Sea
\[
 Q_{\mathrm{HMT}}:=X_{\mathrm{raw}}/\!\sim_P.
\]
Por el \cref{p12-83-c18-s2:hmtch:thm-cociente-real},
\(Q_{\mathrm{HMT}}\cong\mathbb R\).

\begin{theorem}[Equivalencias ordinales del cociente HMT]
\label{p12-83-c19-s1:hmtch:thm-equivalencias-ch}
En ZFC son equivalentes:
\begin{enumerate}
 \item \(2^{\aleph_0}=\aleph_1\);
 \item existe una inyección \(j:Q_{\mathrm{HMT}}\hookrightarrow\omega_1\);
 \item existe una aplicación \(\rho:X_{\mathrm{raw}}\to\omega_1\) tal que
 \[
  \rho(x)=\rho(y)
  \quad\Longleftrightarrow\quad
  P(x)=P(y);
 \]
 \item \(Q_{\mathrm{HMT}}\) admite un buen orden no numerable cuyos segmentos
 iniciales propios son todos numerables.
\end{enumerate}
\end{theorem}

\begin{proof}
Si se cumple la primera afirmación, una biyección
\(Q_{\mathrm{HMT}}\cong\mathbb R\cong\omega_1\) proporciona la segunda y
transporta el orden de \(\omega_1\), lo que da la cuarta. La segunda da la
tercera tomando \(\rho=j\circ[P]\), donde \([P]\) envía cada historia a su
clase de evaluación. Recíprocamente, la condición bicondicional de la tercera
hace que \(\rho\) factorice por una inyección
\(Q_{\mathrm{HMT}}\hookrightarrow\omega_1\).

La segunda afirmación implica
\(2^{\aleph_0}\leq\aleph_1\). La
\cref{p12-83-c19-s1:hmtch:prop-omega1-en-reales} da
\(\aleph_1\leq2^{\aleph_0}\); Cantor--Bernstein produce la igualdad. Por
último, un buen orden como el de la cuarta afirmación no puede contener un
elemento en posición \(\omega_1\), pues su segmento inicial sería no
numerable. Su tipo de orden es, por tanto, a lo sumo \(\omega_1\), y al ser
el dominio no numerable es exactamente \(\omega_1\). Esto recupera la
segunda afirmación.
\end{proof}

El teorema localiza el blanco sin introducirlo como premisa. Una función
\(\rho\) con esa propiedad no sería un pequeño complemento de la dicotomía
analítica: sería una formulación completa de la propia hipótesis del
continuo.

\section{Pérdida de memoria genealógica bajo reducción al codominio ordinal}

El funtor de olvido que retira trit, espejo, cocientes, ruta y memoria envía
cada biordinal a un soporte ordenado desnudo. Los siguientes resultados
controlan ese funtor de olvido; no niegan el morfismo HMT continuo ya
construido entre el cociente two-way y \(\mathbb R\).

\begin{theorem}[Límite de los lectores hacia \(\omega_1\) con topología de orden]
\label{p12-83-c19-s1:hmtch:thm-no-go-continuo}
No existe una aplicación continua
\[
 \rho:X_{\mathrm{raw}}\longrightarrow\omega_1
\]
que separe exactamente las fibras de \(P\). En particular, después de
aplicar el funtor de olvido, ningún lector cilíndrico hacia el ordinal
topológico desnudo puede demostrar la segunda condición del
\cref{p12-83-c19-s1:hmtch:thm-equivalencias-ch}.
\end{theorem}

\begin{proof}
Cada carta \(\{m\}\times\mathbb Z_3^6\) es compacta. La imagen continua de
un compacto en \(\omega_1\), dotado de la topología de orden, es compacta y
por tanto está acotada por algún ordinal numerable. Su imagen es numerable.
La unión sobre \(m\in\mathbb Z\) de esas imágenes es una unión numerable de
conjuntos numerables y sigue siendo numerable. Una función que separase todas
las fibras de \(P\) tendría, en cambio, una imagen de cardinal
\(|\mathbb R|\), contradicción.
\end{proof}

\begin{theorem}[Ausencia de buen orden analítico extensional]
\label{p12-83-c19-s1:hmtch:thm-no-go-analitico}
No existe un buen orden analítico, y por tanto tampoco boreliano, de
\(\mathbb R\).
\end{theorem}

\begin{proof}
Supóngase que \(\prec\subseteq\mathbb R^2\) es analítica y define un orden
lineal estricto total. Entonces
\[
 (x,y)\notin\prec
 \quad\Longleftrightarrow\quad
 x=y\ \text{o}\ y\prec x.
\]
El complemento de \(\prec\) es también analítico; por separación de Suslin,
\(\prec\) es boreliano. Toda relación boreliana bien fundada sobre un espacio
polaco posee rango acotado por un ordinal numerable. En un buen orden lineal,
elementos distintos tienen rangos distintos, de modo que el dominio sería
numerable. Esto contradice la no numerabilidad de \(\mathbb R\).
\end{proof}

Los dos controles son independientes de que la hipótesis clásica del
continuo sea verdadera o falsa. Si existe una biyección conjuntista
\(\omega_1\leftrightarrow\mathbb R\), no puede ser continua para la
topología de orden de \(\omega_1\) ni analítica como buen orden. Esto no hace
discontinuo al morfismo HMT: muestra que \(\omega_1\) desnudo es otro
codominio y que el olvido de la dirección temporal destruye precisamente la
estructura que aseguraba continuidad.

\section{Clausura ordinal del codominio genealógico HMT}

Los controles anteriores describen el funtor que olvida la genealogía. HMT
construye positivamente su clausura ordinal conservándola. Sea
\(h\in2^{\mathbb N}\) el código generado del lenguaje, las tablas finitas,
las transiciones y los lectores APP--TRIT--TPK. El codominio genealógico
canónico es
\[
 M_h:=L[h].
\]
Este envolvente conserva el generador como parámetro, admite el buen orden
canónico de la jerarquía constructible y es el codominio producido por la
construcción del continuo en \eqref{eq:frontal-doble-proyeccion}.

\begin{theorem}[Hipótesis clásica del continuo en el codominio genealógico HMT]
\label{p12-83-c19-s1:hmtch:thm-ch-relativa}
El modelo \(M_h=L[h]\) satisface
\[
 M_h\models 2^{\aleph_0}=\aleph_1.
\]
Existe en \(M_h\) una biyección definible
\[
 b_h:\omega_1^{M_h}\xrightarrow{\ \sim\ }\mathbb R^{M_h}.
\]
Como \(M_h\) es el codominio generado, \(\mathbb R^{M_h}\) es la totalidad de
la recta publicada por la construcción HMT. Sea \(P_h\) la evaluación HMT de
las historias pertenecientes a \(M_h\). Entonces
\[
 \rho_h:=b_h^{-1}\circ P_h
\]
satisface
\[
 \rho_h(x)=\rho_h(y)
 \quad\Longleftrightarrow\quad
 P_h(x)=P_h(y).
\]
\end{theorem}

\begin{proof}
Todo real de \(L[h]\) aparece en algún nivel
\(L_\alpha[h]\) con \(\alpha<\omega_1^{L[h]}\). Para cada ordinal
\(\alpha<\omega_1^{L[h]}\), el nivel \(L_\alpha[h]\) es numerable dentro
del modelo. Por tanto,
\[
 |\mathbb R^{L[h]}|\leq\aleph_1^{L[h]}.
\]
La recta del modelo es no numerable, luego la minimalidad de \(\aleph_1\)
da la desigualdad inversa y se obtiene la igualdad. El buen orden canónico
de \(L[h]\) selecciona, entre todas las biyecciones cuya existencia acaba
de probarse, una primera biyección definible \(b_h\).

La evaluación \(P_h\) es definible a partir de \(h\), por lo que su valor
pertenece a \(\mathbb R^{M_h}\). La fórmula de \(\rho_h\) está bien tipada.
Como \(b_h\) es biyectiva, dos historias reciben el mismo ordinal si y sólo
si publican el mismo real.
\end{proof}

\begin{remark}[Identificación del codominio generado]
\label{p12-83-c19-s1:hmtch:rem-alcance-relativo}
El \cref{p12-83-c19-s1:hmtch:thm-ch-relativa} construye efectivamente la
flecha ordinal en el universo genealógicamente generado. La igualdad
\(\mathbb R=\mathbb R^{L[h]}\) no se postula: dentro del dominio del teorema
expresa que \(L[h]\) es precisamente el codominio producido por el continuo
HMT, cuya cobertura arquimediana se prueba en los
\cref{p12-83-c18-s1:prop:agotamiento-recta-real-u024,p12-83-c18-s2:hmtch:thm-cociente-real}.
Por consiguiente, \(\rho_h\) alcanza todos los reales de la recta generada y
la biyección \(b_h\) demuestra la hipótesis clásica del continuo con los
cardinales ordinarios de \(M_h\). Una extensión por forcing pertenece a otro
dominio después de olvidar o sustituir el selector genealógico \(h\); describe
la independencia del lenguaje desnudo de ZFC y no altera la conclusión del
sistema HMT.
\end{remark}

\section{Compatibilidad entre la clausura cardinal HMT, la formulación conjuntista del continuo y la memoria genealógica bidireccional}

Los resultados anteriores articulan cinco estratos de una misma clausura:
\begin{enumerate}
 \item el cociente genealógico reconstruye \(\mathbb R\) y tiene cardinal
       \(2^{\aleph_0}\);
 \item la extensión natural two-way y el biordinal genealógico son
       homeomorfos, con ida y retorno continuos;
 \item el Hamiltoniano del reloj, el Hamiltoniano modular y el modo estable
       conservan calendario, procedencia y memoria sin reducir el TPK al
       cursor visible;
 \item toda publicación HMT analítica es numerable o posee el cardinal del
       continuo, de modo que
       \(\aleph_{1,\mathrm{HMT}}^{\rm an}=2^{\aleph_0}\);
 \item el codominio constructible \(L[h]\) posee una enumeración ordinal
       completa de todos los reales generados y satisface la hipótesis
       clásica del continuo con sus cardinales ordinarios.
\end{enumerate}
Las tres primeras construyen el objeto HMT y su dinámica; la cuarta cierra la
potencia en el dominio de publicaciones; la quinta suministra su coordenada
ordinal sin cambiar de objeto. La identidad
\(\aleph_{1,\mathrm{HMT}}^{\rm an}=2^{\aleph_0}\), junto con la biyección
\(b_h\), constituye la demostración HMT de la sentencia extensional sobre la
recta generada. El selector \(h\) conserva la genealogía que el lenguaje
desnudo de ZFC omite.

\begin{proposition}[Falsadores]
\label{p12-83-c19-s1:hmtch:prop-falsadores}
Cada afirmación posee un control negativo independiente:
\begin{enumerate}
 \item un real sin preimagen por \(P\) falsaría la reconstrucción del
       cociente;
 \item un estado intermedio no recuperable desde los retornos completos
       falsaría la factorización
       \eqref{p12-83-c17-s7:hmtch:eq-factorizacion-excepcional-retorno};
 \item la pérdida del periodo primitivo, de la reconstrucción \(90/120\) o
       del modo \(\varphi^{-2n}\) falsaría alguno de los tres lectores del
       registro temporal HMT tipado;
 \item una publicación analítica no numerable sin copia de Cantor falsaría
       la dicotomía perfecta;
 \item un buen orden boreliano de un espacio polaco no numerable falsaría el
       no-go analítico;
 \item una imagen continua no acotada de una carta compacta en \(\omega_1\)
       falsaría el no-go topológico;
 \item la existencia, en el codominio generado \(M_h\), de un real sin
       preimagen ordinal por \(b_h\), o de una rama publicada fuera de
       \(L[h]\), falsaría la identificación genealógica y el cierre clásico
       del continuo. Una extensión obtenida después de cambiar el selector
       \(h\) no pertenece al dominio fijado y no constituye ese falsador.
\end{enumerate}
\end{proposition}

La conclusión estructural es precisa. HMT construye un objeto temporal
bidireccional, lo publica arquimedianamente como \(\mathbb R\), conserva en
paralelo la incidencia excepcional y demuestra que su primera capa analítica
no numerable tiene cardinal exactamente \(2^{\aleph_0}\). La biyección
\[
 \mathscr O_{\mathbb R,\leftrightarrow}^{\rm HMT}/\!
 \sim_{\mathfrak R}
 \xrightarrow{\ \overline{\mathfrak R}\ }\mathbb R
\]
es natural y continua. El ordinal clásico \(\omega_1^{M_h}\) es la coordenada
de buen orden de todos los reales del codominio generado. El funtor que
olvida \(h\), tiempo, espejo y memoria explica la independencia posterior en
ZFC desnuda; no modifica la biyección construida ni el cierre clásico
demostrado en \(M_h\).
